\section*{Evaluasi}
\addcontentsline{toc}{section}{Evaluasi}

\begin{enumerate}[label=\arabic*.,leftmargin=*]
  \item Tuliskan definisi formal tata bahasa dan jelaskan setiap
        komponennya.
  \item Klasifikasikan tata bahasa \(S \to aSb \mid \varepsilon\) dan
        jelaskan alasannya.
  \item Buktikan bahwa \(\{a^n b^n \mid n \ge 0\}\) bukan bahasa
        regular.
  \item Jelaskan mengapa \(\{a^n b^n c^n \mid n \ge 1\}\) bukan bahasa
        bebas konteks.
  \item Gambarkan PDA untuk \(\{a^n b^{2n} \mid n \ge 0\}\).
  \item Bandingkan daya ekspresif DFA, PDA, dan mesin Turing.
  \item Mengapa lexer cukup memakai otomata hingga, sedangkan parser
        memerlukan otomata pushdown?
  \item Berikan satu contoh kesalahan yang tidak dapat dideteksi oleh
        tata bahasa bebas konteks.
\end{enumerate}

\subsection*{Kunci dan Rubrik Singkat}
\begin{itemize}[leftmargin=*]
  \item Soal 1--2: ketepatan definisi dan klasifikasi (70\%),
        penjelasan (30\%).
  \item Soal 3--4: struktur pembuktian (60\%), ketepatan argumen
        (40\%).
  \item Soal 5: kebenaran transisi dan stack (80\%), kelengkapan
        (20\%).
  \item Soal 6--8: ketepatan konsep (60\%), contoh (40\%).
\end{itemize}
